% TOP_PROBLEM: {"id": 3247, "title": "Weighted colouring of hexagonal graphs.", "category_id": 3, "category": {"id": 3, "name": "graph_theory", "display_name": "Graph Theory", "description": "Problems involving graphs, networks, and their properties.", "slug": "graph-theory", "order_index": 3, "created_at": "2026-07-31T15:26:25.671Z"}, "status": "open", "problem_number": "OPG-47511", "set_id": 12}
% FIRST_POSTED: 2026-10-01
% ENTRY_KIND: statement_only
% UnsolvedMath ID 3247; source catalogue code OPG-47511
% !TeX program = lualatex
% Definition and sources only; this file is not an LLM solution attempt.
\documentclass[11pt,a4paper]{article}
\usepackage[margin=25mm]{geometry}
\usepackage{fontspec}
\setmainfont{lmroman10-regular.otf}[BoldFont=lmroman10-bold.otf,
 ItalicFont=lmroman10-italic.otf,BoldItalicFont=lmroman10-bolditalic.otf]
\usepackage{amsmath,amssymb,mathtools}
\usepackage{unicode-math}
\setmathfont{latinmodern-math.otf}
\AtBeginDocument{\let\setminus\smallsetminus}
\usepackage{xurl,hyperref}
\hypersetup{colorlinks=true,urlcolor=blue}
\setcounter{secnumdepth}{0}
\setlength{\parindent}{0pt}
\setlength{\parskip}{5pt}
\setlength{\emergencystretch}{3em}
\begin{document}
\section*{Weighted colouring of hexagonal graphs.}\label{3247}
{\small UnsolvedMath ID 3247 · OPG-47511 · Graph Theory · OpenGarden}
\subsection{Definitions and mathematical statement}
Let $e_1=(1,0)$ and $e_2=(1/2,\sqrt3/2)$. The triangular lattice has vertices $ae_1+be_2$ for $a,b\in\mathbb Z$ and joins pairs at Euclidean distance one. A hexagonal graph in this problem is a finite induced subgraph of that lattice: once its vertex set is chosen, every lattice edge between selected vertices is retained. This terminology does not mean a subgraph of the degree-three honeycomb lattice.

Give each vertex a positive integer demand $p(v)$. A multicoloring with $k$ colors assigns a set $C(v)\subseteq\{1,\ldots,k\}$ of exactly $p(v)$ labels to each vertex, with $C(u)\cap C(v)=\varnothing$ whenever $uv$ is an edge. Let $\chi(G,p)$ be the smallest such $k$. Define the weighted clique number by
\[
 \omega(G,p)=\max\left\{\sum_{v\in Q}p(v):Q\subseteq V(G)
                          \text{ is a clique}\right\},
\]
where a clique has all its distinct pairs adjacent; singleton cliques are included.

The conjecture is that there exists an absolute constant $c\ge0$ such that every hexagonal graph and every positive integer demand function satisfy
\[
                         \chi(G,p)\le\tfrac98\omega(G,p)+c.
\]
The additive constant is independent of both graph size and demand magnitudes. Demands specify distinct colors within each vertex, rather than weights in an ordinary one-color assignment. No triangle-free assumption is made in this main conjecture, and triangles contribute their full combined demand to the clique lower bound.
\subsection{Short English statement}
Can every demanded multicoloring of an induced triangular-lattice graph be achieved within nine-eighths of its weighted clique bound plus a universal constant?
\subsection{Sources}
\begin{itemize}
\item[S1] ULAM.AI and the UnsolvedMath Contributors, supplied problems.json, record 3247 (OPG-47511), OpenGarden. Catalogue identity and source transcription; CC BY 4.0.
\url{https://huggingface.co/datasets/ulamai/UnsolvedMath}.
\item[S2] Open Problem Garden, Weighted colouring of hexagonal graphs.. Original problem and discussion.
\url{http://www.openproblemgarden.org/op/weighted_colouring_of_hexagonal_graphs}.
\end{itemize}
\end{document}
